\chapter{绪论——从哲学本体论到工程本体论}
\chapterart{ch01}{从哲学之问到工程之答：本体论的两种形态}

\begin{introbox}
\textbf{【导读】}
为什么一本讲工程的书要从哲学谈起？因为「本体论」这个词承载着两千多年的哲学追问，
而我们要做的，是把它改造成可以写进代码、可以被推理器检验的工程工具。
本章从制造车间的三个典型教学困境出发，先看清企业知识「信息都在、知识不通」的高熵现状；
然后回到哲学源头，用克制的篇幅交代亚里士多德范畴论以来的思想脉络，
理解哲学追问「无法收敛、不可计算」的特征；
接着进入本书的主战场——工程转向：逐词精读 Gruber、Borst 与 Studer 的经典定义，
掂量概念化、显式、形式化、共享四个关键词的工程分量；
再用「知识熵减」这个统一视角串起全书方法，并把本体与数据库 Schema、分类法、
叙词表、UML 类图、知识图谱等相邻概念逐一辨析清楚；
最后回答 AI 时代的核心问题——大模型为什么需要受控事实来源与 Semantica 语义权威层，
并给出全书的知识地图与三条学习路径。
读完本章，你应当能用一句话向同事解释：本体不是哲学玄谈，而是一套让机器按显式边界处理业务语义的基础设施。
\end{introbox}

\chapterbinding{semantica.chapter\_packages.vol1.ch01}{partial}{blocked}
{章合同、CQ、两份概念案例与全书路线图已经迁入；包可被原生 registry 装载。}
{ontology、SHACL、SPARQL 以及可执行 conceptualization 边界报告均缺失；
当前场景与 oracle 为 absent，不能把阅读练习写成机器验收。}

\section{从一个车间的清晨说起}

本书的主角是「知识」——不是泛泛而谈的知识，而是让一家企业能够正常运转的那些具体判断：
哪台设备能干哪种活、哪道工序必须排在哪道之前、哪个参数在什么条件下要调整。
在给出任何定义之前，我们先到一个机加工车间里待一个上午，看三个每天都在发生的场景。
读者会发现，它们表面上是三个不同部门的烦恼，病根却是同一个。

\subsection{场景一：排产决策——知识在谁的脑子里}

早上七点半，调度员小王坐到工位前，面对三块屏幕：
制造执行系统（MES，Manufacturing Execution System）里挂着今天要排的四十张工单，
Excel 台账里躺着两百台设备的型号与参数，
企业资源计划系统（ERP）里还有物料齐套情况和每张订单的交期。
屏幕上的信息不可谓不全。但真正决定「怎么排」的知识——
「3 号车床不能干钛合金」「热处理必须排在焊接后面」
「A 区那台加工中心的精度最近不稳，精密件先别往上放」——
不在任何系统里，在小王和几位老师傅的脑子里。

把这些知识摊开看，大致有三类。
第一类是\textbf{设备能力知识}：每台设备能加工什么材料、什么尺寸范围、达到什么精度等级。
台账里有功率、行程这些「写得出来的」参数，
但「3 号车床干不了钛合金」这种结论性的能力判断，恰恰不在台账的任何字段里——
它是某次崩刃事故之后老师傅总结出来的，原因涉及冷却方式与转速区间的组合，没人把它写下来。
第二类是\textbf{工艺约束知识}：工序之间的先后关系、设备与工装的搭配禁忌、
不同材料混线生产时的清洁要求。这些约束一部分写在工艺文件里，
另一部分以「车间惯例」的形式存在——新来的工艺员往往要靠犯错来学习它们。
第三类是\textbf{动态状态知识}：哪台设备最近精度漂移、哪位操作工今天请假、
哪批毛坯料的硬度偏高。这类知识的保鲜期只有几天，几乎完全依赖口头传递。

去年小王休年假，顶班的小李按台账把一批钛合金件排给了 3 号车床——
台账上功率与扭矩都够。结果干了不到一个小时就崩刃停机，
零件报废两件，重排计划波及后面三天的交付。
事后复盘，谁都没有过错：台账没有错，工单没有错，小李的推理也没有错，
错在那条决定性的知识只存在于不在场者的脑子里。
这就是排产场景的真实困境：\textbf{数据可以买服务器来存，
这类决定成败的判断却只能靠老师傅的在场来「存」}。

更值得玩味的是信息化的悖论。这家工厂先后上了 MES、ERP 和高级排产系统（APS），
数据一年比一年多，可自动排产方案每次生成出来，
老师傅扫一眼就能挑出几处「这么排要出事」的地方——系统不知道的约束，全在人的脑子里。
于是 APS 沦为草稿生成器，最终版本仍然靠人手工调整。
问题不在算法不够聪明，而在\textbf{算法能拿到的知识太少}：
约束没有被显式地、形式化地交给机器，机器自然排不出可执行的计划。

\subsection{场景二：新人培训——老师傅的经验为什么传不下去}

王工在热处理车间干了三十二年，明年退休。
厂里给他安排了两件事：带三个徒弟，以及「把经验写下来」。
三个月后王工交来一份四十多页的《热处理操作要点》，里面写满了这样的句子：
「视材料状态适当延长保温时间」「炉温接近上限时应注意观察」「装炉量大时酌情调整」。
徒弟拿着文档来问：什么叫「适当」？多少算「接近」？「酌情」是看什么情？
王工想了想，说：这说不清楚，得看料的批次，看炉子的脾气，跟着我看几年就懂了。

王工没有藏私，他遇到的是知识管理中的经典难题：\textbf{隐性知识}（tacit knowledge）。
哲学家波兰尼（Michael Polanyi）有一句常被引用的话：「我们知道的，多于我们能说出的。」
骑自行车的人说不清自己如何保持平衡，老师傅也说不清「炉子的脾气」具体指哪些变量。
但如果就此把所有经验都归入「不可言传」，企业就只能接受知识随人退休而流失。
仔细分析会发现，老师傅的经验其实分两类：
一类是\textbf{技能型}的，如听切削声音判断刀具磨损——这类知识与身体感知绑定，
确实难以形式化，本书不处理它们；
另一类是\textbf{规则型}的，如「这批料硬度偏高，转速要降一档」——
它完全可以被形式化，只是从来没有人逼问出它的精确条件：
硬度高于多少？降到哪一档？对哪些刀具成立？
工程本体论主攻的正是第二类：\textbf{把「视情况而定」逼问成「视什么情况、定到什么值」}。

企业当然尝试过各种对策。师徒制有效但太慢，覆盖面窄，
徒弟学成与否高度依赖个人悟性，而且人一流动，积累就清零。
文档库与企业 Wiki 看似把知识留住了，实际效果常常令人失望：
检索靠关键词匹配，写文档的人用「淬火」，搜索的人输入「热处理硬化」，就是搜不到；
同一主题的文档有五个版本，谁也不知道哪个最新；
流程早改了，文档没人更新，新人照着过时文档操作反而出错。
还有企业把老师傅的操作录成视频——能看，不能查，更不能让系统拿来做自动校验。
这些对策的共同缺陷是：\textbf{知识没有结构}。
一段文字、一段视频对机器来说只是字节流，机器无法判断其中讲了哪些概念、
概念之间什么关系、新来的一条记录与旧知识是否矛盾。
没有结构，机器就帮不上忙；机器帮不上忙，知识传承就只能继续依赖人盯人。

\subsection{场景三：质量审计——断裂的追溯链}

第三个场景发生在质量部。某天上午，客户投诉上月交付的一批法兰盘硬度抽检不合格，
要求工厂四十八小时内给出解释。质量工程师老郑需要回答三个问题：
这批零件是哪台炉子、按什么参数处理的？
同一炉次的其他批次流向了哪些客户，是否需要追加排查？
当时的工艺参数为什么这样设置，是谁、依据什么批准的？

老郑的追溯之旅是这样的：批次号在 ERP 里查到对应的销售订单与生产工单；
拿工单号到 MES 里查到热处理工序的报工记录，记录里的设备字段写的是「资产编号」；
而热处理车间那台炉子的控制系统是单机软件，导出的记录用的是「设备编号」，
两套编号规则不同，对应关系只存在于一张多年前的 Excel 对照表和老郑的记忆里；
保温曲线存在仪表厂商的私有格式文件中，要找厂商的工具才能打开；
最关键的「为什么调整保温时间」，最后在王工的纸质交接本上找到一行字：
「该批料含碳量偏上限，延保三十分钟，已电话告知质检」。
三天后，老郑拼出了一条「大概率正确」的解释链。
客户的审计师看完材料只问了一句：这两个编号是同一台设备，有系统证据吗？
老郑答不上来——那一环靠的是「我记得」。

这个场景暴露的问题与前两个不同：信息其实都被记录了，
但它们\textbf{以互相无法对接的方式}被记录在不同载体里。
「设备」这个概念在 ERP 里叫资产，在 MES 里叫设备，在车间单机系统里叫机台，
三套标识符、三种字段结构，没有任何一处声明「它们指的是同一类东西、这两条记录指的是同一台」。
审计要求的是\textbf{证据链}——每一环都可验证、可复查；
而系统能提供的只是\textbf{数据片段}，环与环之间的连接靠人的记忆与口头解释。
在普通机加工行业，这意味着多花三天和一次尴尬；
在航空、医疗器械、核电等强监管行业，可追溯性（traceability）是市场准入条件，
追溯链断裂直接意味着资质风险。

\subsection{诊断：信息都在，知识不通}

三个场景来自三个部门，表面上是三种烦恼：排产怕人不在，培训怕人退休，审计怕链条断。
把它们放在一起看，病根是同一个，可以概括成一句话：\textbf{信息都在，知识不通}。
进一步解剖，「不通」有三个具体病灶：

\begin{itemize}
\item \textbf{分散}（fragmentation）：数据散落在 MES、ERP、Excel、单机系统、纸质记录中，
      每个系统是一座孤岛，岛与岛之间没有共同的语义层；
\item \textbf{隐性}（tacitness）：真正决定行动的判断——能力边界、工艺禁忌、调整理由——
      存在于个人经验中，系统里只有这些判断留下的「数据残影」；
\item \textbf{歧义}（ambiguity）：同一个概念在不同部门有不同名称（同义异名），
      同一个名称在不同系统里指不同东西（同名异义），
      没有任何机制保证大家说的是同一件事。
\end{itemize}

\begin{center}
\begin{tabularx}{\linewidth}{@{}lXX@{}}
\toprule
\textbf{场景} & \textbf{表面症状} & \textbf{深层病灶} \\
\midrule
排产决策 & 顶班者按台账排产导致崩刃；APS 方案无法直接执行 & 能力与约束知识隐性化，机器拿不到决定性规则 \\
新人培训 & 经验文档充满「适当」「酌情」；检索不到、版本混乱 & 规则型经验未被精确化，文档无结构、机器不可用 \\
质量审计 & 追溯耗时三天，关键环节靠记忆作证 & 概念与标识分散且歧义，系统之间没有共同语义层 \\
\bottomrule
\end{tabularx}
\end{center}

这是几乎所有制造企业——推而广之，几乎所有知识密集型组织——的知识现状。
用物理学的比喻说，企业知识处于「高熵」状态：混乱、隐性、无法被机器利用。
工程本体论要做的事，本质上就是给知识做\textbf{熵减}：
把隐性的变成显式的（老师傅的经验写成形式化规则），
把分散的变成统一的（全企业共享一套概念定义），
把模糊的变成可推理的（机器能基于规则自动检查、自动推导）。
这本书讲的就是这件事的原理、方法和工程实践。

在动手之前，还需要先回答一个绕不开的问题：
「本体论」这个听起来像哲学课的词，怎么会跟车间里的这些麻烦扯上关系？
这要从两千多年前说起——放心，我们只在哲学里停留一节的篇幅，
很快就会回到工程师的世界。

\section{本体论的哲学源流}

「本体论」译自英文 Ontology，词根来自希腊语：on（存在者）加上 logos（学说），
合起来就是「关于存在的学说」。
拉丁语合成词 ontologia 直到十七世纪初才在哲学著作中出现，
但它所指的那门学问，从古希腊就已经开始。
本节用最克制的篇幅交代这条思想脉络——不是为了学术礼仪，
而是因为工程本体论的几件核心工具（分类层次、定义方法、同一性判据）都是从这里继承的，
知道源头，用起来才有分寸。

\subsection{亚里士多德：范畴与属加种差}

公认的起点是亚里士多德。他在《形而上学》中提出要研究「存在之为存在」
（being qua being）：不研究某一类存在者——那是各门具体科学的事——
而研究一切存在者作为存在者的普遍规定。
在《范畴篇》中，他给出了著名的十范畴：实体、数量、性质、关系、地点、时间、
姿态、状况、动作、承受。
任何一个可以被谈论的东西，都落入这十个最高的「类」之一。
这件事的开创性在于：\textbf{它是人类第一次对「能够谈论的一切」做顶层分类}。
今天任何一个本体工程项目的第一步——「这个领域里有哪几大类东西？」——
做的仍然是同一件事，只是范围从「一切存在」收缩到了某个业务领域。

亚里士多德还留下了一件更趁手的工具：\textbf{属加种差}（genus plus differentia）的定义法。
要定义一个概念，先指出它的上位类（属），再给出它区别于同属其他成员的特征（种差）：
「人是理性的动物」——动物是属，理性是种差。
两千多年后，OWL 语言里定义一个类的标准句式仍然是这个结构：
「高端设备 \(\equiv\) 设备 \(\sqcap\) 若干特征约束」——属是「设备」，种差是后面的约束。
第 2 章将看到这种定义的严格写法。

同样在《范畴篇》里，亚里士多德区分了「第一实体」（个体，如苏格拉底这个人）
与「第二实体」（种与属，如「人」「动物」这些类）。
个体与类的两层结构，正是今天知识库中实例层与概念层的远祖——
第 2 章的 ABox 与 TBox 将再次出现这对区分。

\subsection{波菲利树：最早的概念层次图}

公元三世纪，新柏拉图主义者波菲利（Porphyry）为《范畴篇》写了一篇导论，
后人根据其中的论述画出了「波菲利树」（Tree of Porphyry）——
思想史上第一张广为流传的概念层次图：
实体之下分有形体与无形体，有形体之下分有生命与无生命，
有生命之下分动物与植物，动物之下分理性与非理性，
理性的动物即人，人之下是苏格拉底、柏拉图这些个体。

这棵树的形态特征——自顶向下、逐层二分、中间是类、叶子是个体——被后世无数次复制：
林奈的生物分类体系、图书馆的图书分类法、企业的设备分类树、
面向对象编程的类继承体系，都是它的后裔。
但它也埋下了一个一直延续到今天的隐患：\textbf{单一维度的树撑不住多视角的世界}。
一台激光切割机，按动力源分是电气设备，按工艺分是切割设备，
按精度等级分是精密设备——三种分法都有道理，强行塞进一棵树就会互相打架。
如何处理多重分类，是第 3 章本体建模方法论中的重要议题。

\subsection{共相之争与近代之后}

中世纪经院哲学围绕一个问题争论了数百年：\textbf{共相}（universals），
即「人」「白色」这类普遍概念，究竟真实存在吗？
实在论者认为共相先于并独立于个体事物而存在；
唯名论者认为真实存在的只有个体，共相不过是名称；
概念论者居中：共相作为概念存在于心智之中。
这场争论本身不需要工程师选边站，但值得注意的是：
当我们在本体中写下 \texttt{CNCMachine} 这个类时，实际上已经采取了一种温和的概念论立场——
类是为了组织知识而约定的概念，我们不追问它的形而上学地位，
判据从「真不真」悄悄变成了「好不好用」。这个转变正是下一节工程转向的种子。

近代以后，哲学的重心从本体论转向认识论；
康德进一步裁定「物自体」不可知，传统本体论的雄心受到重挫。
二十世纪，分析哲学用现代逻辑重新表述存在问题，
蒯因（W. V. O. Quine）给出了一个对工程师极有启发的说法：
「存在就是成为约束变元的值」——一个理论承诺哪些东西存在，
就看它的量词变元需要在哪些东西上取值。
把「理论」换成「信息系统」，这句话几乎就是工程本体论的纲领：
\textbf{一个系统承诺哪些东西存在，就看它的本体里声明了哪些类与个体}。

哲学传统还锤炼了一对工程上极易混淆的关系：「是一种」与「是其部分」。
逻辑学家与哲学家为后者发展出专门的理论——部分论（mereology），
仔细区分了「发动机是汽车的部件」与「轿车是汽车的一种」这两类形似而实异的陈述。
这个区分将在本书中反复出现：1.5 节的叙词表正因混淆二者而无法支撑推理，
第 3 章的建模方法论也会把「误将部分当子类」列为高频建模错误。

回望这两千多年，哲学本体论的问题——世界的本源、事物的同一性、共相的地位——
有一个共同特征：\textbf{无法收敛，不可计算}。
每一代哲学家都给出新的回答，没有一个回答能够终结讨论；
这些问题也无法交给任何程序去判定。
包内 \texttt{philosophy-to-engineering} 资产的第一节把这条线索压缩成几行注释：

\begin{codebox}
\codeline{\textcolor{cmtgray}{\#\ ==============================}}
\codeline{\textcolor{cmtgray}{\#\ 1.\ 哲学本体论的追问}}
\codeline{\textcolor{cmtgray}{\#\ ==============================}}
\codeline{\textcolor{cmtgray}{\#\ 哲学本体论（首字母大写的\ Ontology）研究"存在之为存在"：}}
\codeline{\textcolor{cmtgray}{\#\ \ \ 世界的本源是什么？}}
\codeline{\textcolor{cmtgray}{\#\ \ \ 一个事物凭什么是"它自己"？}}
\codeline{\textcolor{cmtgray}{\#\ \ \ 共相（普遍概念）是否真实存在？}}
\codeblank{}
\codeline{\textcolor{cmtgray}{\#\ 亚里士多德的范畴论给出了最早的"顶层分类"：}}
\codeline{\textcolor{cmtgray}{\#\ \ \ 实体（Substance）、数量（Quantity）、性质（Quality）、}}
\codeline{\textcolor{cmtgray}{\#\ \ \ 关系（Relation）、地点（Place）、时间（Time）……}}
\codeblank{}
\codeline{\textcolor{cmtgray}{\#\ 特征：终极追问无法收敛，知识永远处于开放的不确定性中}}
\codeline{\textcolor{cmtgray}{\#\ ——\ 这类问题"无法落地、不可计算"，不是工程本体论的研究对象}}
\codeblank{}
\end{codebox}
\color{black}

\assetsource{ch01}{philosophy-to-engineering}

请注意注释最后一行的措辞：「无法落地、不可计算」不是对哲学的贬低，而是一次\textbf{划界}。
工程本体论从哲学那里接手的不是答案，而是问题清单与概念工具：
顶层分类的思想、属加种差的定义术、对同一性判据的敏感。
至于那些终极追问本身，工程师礼貌地把它们留在哲学系。
本书对哲学源流的交代到此为止；从下一节起，
「本体」一词将默认指工程意义上的本体——
若读者对哲学脉络意犹未尽，任何一部严肃的西方哲学史都比本书讲得好。

\section{工程转向：从 Ontology 到 ontologies}

\subsection{转向的背景：知识为什么需要共享}

二十世纪七八十年代，人工智能的主流形态是\textbf{专家系统}（expert system）：
把领域专家的判断规则编码进程序，让机器在窄领域内模仿专家决策。
专家系统在医疗诊断、设备配置等领域取得过商业成功，
但很快撞上了著名的\textbf{知识获取瓶颈}（knowledge acquisition bottleneck）：
把专家知识变成规则的过程极其昂贵缓慢；
更糟的是，每个系统的知识库都是从零开始、用各自的私有结构编码的——
A 系统辛苦积累的设备知识，B 系统一行都用不上。
知识无法共享与复用，意味着每个项目都在重复支付最贵的那笔成本。

九十年代初，知识工程界发起了以知识共享为目标的若干研究计划，核心思路是：
把知识库拆成两层——随任务而变的部分，
和\textbf{相对稳定、可以在系统之间共享的概念体系}。
后者需要一个名字。研究者们从哲学那里借来了 ontology 这个词，
用它指称「对某个领域里存在哪些概念、概念之间有哪些关系的明确约定」。
1993 年，Gruber 给出了被引用最多的定义：

\begin{noteblock}
\textbf{Gruber（1993）}：An ontology is an explicit specification of a conceptualization.
——本体是概念化的显式规范。
\end{noteblock}

随后 Borst（1997）补充了两个限定：本体应当是\textbf{形式化}的规范，
且概念化应当是\textbf{共享}的；
Studer 等人（1998）把两者综合成今天教科书通用的版本：

\begin{noteblock}
\textbf{Studer 等（1998）}：本体是\textbf{共享概念化}的\textbf{形式化}、\textbf{显式}规范
（a formal, explicit specification of a shared conceptualization）。
\end{noteblock}

这个定义浓缩了四个关键词：概念化、显式、形式化、共享。
它们不是修辞排比，每一个都对应一条工程要求、划掉一类「不是本体」的东西。
下面逐词精读，每个词配一个制造业的例子。

\subsection{关键词一：概念化——选定看世界的侧面}

\textbf{概念化}（conceptualization）指对世界某个侧面的抽象模型：
从纷繁的现实中，选出我们关心的对象类型、关系与约束，忽略其余。
关键在「侧面」二字——概念化必然是\textbf{有选择的、视角依赖的}。
同一台车床，在排产视角下是一个「产能单元」，关心的是它能加工什么、节拍多少、何时空闲；
在维修视角下是一个「部件组合体」，关心的是主轴、导轨、丝杠的状态与故障模式；
在财务视角下是一项「固定资产」，关心的是原值、折旧与残值。
三种概念化都正确，但服务于不同任务。
所以本体工程的第一课就是：\textbf{不存在「把企业的一切都建进去」的万能本体}，
只存在为明确任务服务的概念化——这也是第 3 章「能力问题」方法的理论根据。

具体地说，一次概念化要回答三个问题：
这个领域里有哪些\textbf{概念}（类）？——设备、工序、物料、工单、班次。
概念之间有哪些\textbf{关系}（属性）？——设备「执行」工序、工序「消耗」物料、工序「先于」工序。
哪些\textbf{约束}必须成立？——一台设备同一时刻至多执行一道工序。
注意，此时概念化可以只存在于建模者的头脑或白板上——它还不是本体，只是本体的原料。

\subsection{关键词二：显式——把约定摆到明处}

\textbf{显式}（explicit）要求概念、关系与约束被明确陈述出来，
而不是隐含在某个人的头脑、某段程序的逻辑或某份没人读的文档的字里行间。
对照场景一：「3 号车床不能干钛合金」长期处于隐式状态——
它真实存在并支配着排产决策，但任何系统都看不见它。
显式化这条知识的正确姿势，不是简单给 3 号车床贴一个「禁钛」标签，
而是把背后的一般性条件挖出来并写明：
加工钛合金要求机床具备高压冷却能力、转速可控制在低速大扭矩区间、整机刚性达标；
然后分别声明 3 号车床的实际配置。
这样「3 号车床不能干钛合金」就从一条孤立的禁令，
变成了一条\textbf{可以由一般规则推导出来的结论}——
并且当 3 号车床明年加装高压冷却之后，结论会自动随事实更新。

显式化还有一个常被低估的副产品：\textbf{暴露分歧}。
很多企业第一次组织本体评审会时才发现，
生产部说的「关键设备」是产线瓶颈设备，设备部说的是高维修成本设备，
安全部门说的是有安全风险的设备——
三个部门用同一个词开了十年会，说的从来不是同一批机器。
隐式状态下，分歧被含糊的语言掩盖；显式化逼着大家当场对齐。
这个过程有时充满火药味，但火药味正是价值所在：
\textbf{分歧在建模会议室里爆发，总好过在事故现场爆发}。

\subsection{关键词三：形式化——让机器读得懂}

\textbf{形式化}（formal）要求规范采用机器可读、语义精确、支持自动推理的语言书写，
而不是自然语言。
自然语言的问题不在于啰嗦，而在于\textbf{歧义不可消除}：
「设备需要定期维护」——哪类设备？多久算定期？车床和天车的「定期」一样吗？
「大功率设备优先检修」——多大算大？两台都「大」时谁优先？
人靠语境消化这些歧义，机器不能。
形式化语言用受限的语法和数学定义的语义把歧义压到零，例如：

\begin{noteblock}
\centering
CNCMachine \(\sqsubseteq\) \(\exists\)hasMaintenancePlan.MaintenancePlan
\end{noteblock}

这条公理只有一种读法：「每台数控机床都至少关联一份维护计划」。
推理器可以拿着它自动巡检知识库，找出所有违反约定的数控机床实例。
符号 \(\sqsubseteq\)（蕴含）与 \(\exists\)（存在约束）的严格定义在第 2 章给出，
此处只需体会一点：\textbf{形式化之后，知识从「人读的文本」变成了「机器可执行的约定」}。

需要说明的是，形式化是一个程度问题，存在一条光谱：
最轻的一端是受控词表（只规定术语），向上是分类层次（再加上下位关系），
再向上是带属性与约束的概念模型，最重的一端是带完整逻辑公理的理论。
本体工程界把偏轻的一端称为轻量级本体，偏重的一端称为重量级本体。
形式化程度越高，机器能做的检查与推理越多，建模的学费也越贵——
第 2 章与第 4 章会教你怎么交这笔学费，并讨论如何按任务选择形式化程度。

\subsection{关键词四：共享——一个人的本体不是本体}

\textbf{共享}（shared）指本体表达的是\textbf{相关团体的共识}，而非个人观点。
这是四个关键词里最容易被技术人员忽视、却最决定成败的一个。
道理很直接：本体的全部价值在于让不同的人、不同的系统对同一套概念达成一致；
一份只有建模者自己认账的「本体」，再精致也只是个人笔记。
共享范围划定价值边界：一个团队内共享，得到的是统一的术语表；
全企业共享，得到的是跨系统集成的语义基础设施；
全行业共享，得到的就是行业标准——1.7 节将看到若干实例。

共识不会从天而降，它是\textbf{社会过程}的产物：评审、争论、妥协、签字。
回到「工单」的例子：车间叫派工单，计划部叫生产指令，财务叫成本归集对象。
共享本体并不要求三个部门改口——那既不现实也无必要——
而是显式声明：三个名称是同一概念的\textbf{别名}（synonym），
各系统继续各叫各的，但在语义层都映射到同一个类。
第 3 章的方法论会把「如何组织人达成共识」制度化成具体步骤：
能力问题评审、术语表谈判、专家走查。
这里先记住一句话：\textbf{本体工程一半是逻辑工程，一半是社会工程}。

\subsection{小写复数的 ontologies}

英文文献里有一个意味深长的细节：
哲学的本体论写作大写单数的 \textbf{Ontology}——它是一门学问的专名，
研究「存在」这个唯一的总主题；
工程的本体写作小写可数的 \textbf{ontology}，而且常用复数 \textbf{ontologies}——
它是普通名词，可以说「一个本体」「几个本体」。
一个排产本体、一个维修本体、一个质量本体，可以并存、各为其任。

这个语法细节宣告了定位的根本改变：
本体从「对世界的唯一正确描述」转变为\textbf{工件}（artifact）——
像软件一样被设计、被评审、被打版本号、被重构、甚至被废弃。
没有「终极正确」的本体，只有对当前任务「足够好」的本体。
当然，多本体并存也带来新问题：两个本体对「设备」的定义不一致时怎么对齐？
这是知识融合的议题，第 7 章处理。

\subsection{一个例子看懂转向：忒修斯之船的工程消解}

两种本体论的气质差异，用一个例子可以说透。
哲学家问：一台车床在大修中更换了全部零件，它还是原来那台车床吗？
这是忒修斯之船问题的车间版本，两千年来没有公认答案——它也不需要有。
工程本体论面对同一个问题的做法是\textbf{约定}：
在我们的信息系统里，序列号就是设备同一性的判据；
换了多少零件都不影响身份，序列号变了就视为另一台设备。
同一包资产把这组对照写成如下教学片段：

\begin{codebox}
\codeline{\textcolor{cmtgray}{\#\ ==============================}}
\codeline{\textcolor{cmtgray}{\#\ 3.\ 对照：同一个"设备"概念}}
\codeline{\textcolor{cmtgray}{\#\ ==============================}}
\codeline{\textcolor{cmtgray}{\#\ 哲学视角的追问：}}
\codeline{\textcolor{cmtgray}{\#\ \ \ "设备"作为人造物（Artifact）的本质是什么？}}
\codeline{\textcolor{cmtgray}{\#\ \ \ 一台车床更换全部零件后还是原来那台车床吗？（忒修斯之船）}}
\codeblank{}
\codeline{\textcolor{cmtgray}{\#\ 工程视角的约定：}}
\codeline{Equipment\ \(\sqsubseteq\)\ PhysicalObject\ \ \ \ \ \ \ \ \ \ \ \ \ \ \ \#\ 设备是物理对象}
\codeline{Equipment\ \(\sqsubseteq\)\ \(\exists\)hasSerialNumber.string\ \ \ \ \ \ \#\ 每台设备有序列号（同一性判据）}
\codeline{CNCLathe\ \(\sqsubseteq\)\ CNCMachine\ \(\sqsubseteq\)\ Equipment\ \ \ \ \ \ \ \ \#\ 分类层次}
\codeline{Equipment\ \(\sqcap\)\ Material\ \(\sqsubseteq\)\ \(\bot\)\ \ \ \ \ \ \ \ \ \ \ \ \ \ \ \ \ \#\ 设备与材料不相交}
\codeblank{}
\codeline{\textcolor{cmtgray}{\#\ 工程本体论不回答"本质"问题，}}
\codeline{\textcolor{cmtgray}{\#\ 而是给出可检验、可推理、可共享的形式化约定：}}
\codeline{\textcolor{cmtgray}{\#\ 序列号就是这个信息系统中设备的同一性判据——问题被"约定"解决}}
\codeblank{}
\end{codebox}
\color{black}

\assetsource{ch01}{philosophy-to-engineering}

花一分钟逐行读这四条公理（符号的严格语义在第 2 章，这里先建立直觉）：
第一行说设备是一种物理对象——\(\sqsubseteq\) 读作「是一种」或「蕴含」；
第二行说每台设备至少有一个字符串类型的序列号——\(\exists\) 读作「存在」，
这就是被约定下来的同一性判据；
第三行用两个 \(\sqsubseteq\) 串出分类层次：数控车床是数控机床，数控机床是设备；
第四行说「设备」与「物料」两个类不相交——\(\sqcap\) 是交，\(\bot\) 是空类——
任何东西不能既是设备又是物料，「把 3 号车床当原材料投进熔炉」在这个本体里是非法陈述。

请注意发生了什么：忒修斯之问\textbf{没有被解答，而是被消解了}。
工程师不证明「序列号判据是真理」——换一家企业完全可以约定用资产编码做判据——
工程师只保证三件事：约定是\textbf{显式}的（写成公理而非藏在脑中），
是\textbf{共享}的（全企业系统一体遵守），
是\textbf{可检验}的（机器能依据它判断任何陈述合法与否）。
哲学问题问「什么是真的」，工程问题问「我们约定了什么、约定是否自洽、是否人人可查」。
这就是「工程化」三个字的全部含义：\textbf{把不可判定的追问，转化为可检验的约定}。

\subsection{什么样的本体是好本体}

既然不再以「哲学上正确」为标准，工程本体的好坏靠什么评判？
包内教学资产的最后一节给出五条工程判据：

\begin{codebox}
\codeline{\textcolor{cmtgray}{\#\ ==============================}}
\codeline{\textcolor{cmtgray}{\#\ 5.\ 工程判据}}
\codeline{\textcolor{cmtgray}{\#\ ==============================}}
\codeline{\textcolor{cmtgray}{\#\ 一个"好的"工程本体不以哲学正确性为判据，而以下列工程性质为判据：}}
\codeline{\textcolor{cmtgray}{\#\ \ \ 一致性（Consistency）\ \ ——\ 公理之间无矛盾，推理器可检验}}
\codeline{\textcolor{cmtgray}{\#\ \ \ 完备性（Completeness）\ ——\ 覆盖领域内需要回答的问题（见第3章能力问题）}}
\codeline{\textcolor{cmtgray}{\#\ \ \ 简约性（Minimality）\ \ \ ——\ 不引入与任务无关的概念}}
\codeline{\textcolor{cmtgray}{\#\ \ \ 可扩展性（Extensibility）——\ 新概念可平滑加入而不推翻已有结构}}
\codeline{\textcolor{cmtgray}{\#\ \ \ 共享性（Consensus）\ \ \ \ ——\ 领域专家认可并实际使用}}
\end{codebox}
\color{black}

\assetsource{ch01}{philosophy-to-engineering}

五条判据各自针对一类失败模式。
\textbf{一致性}针对「公理打架」——约束之间存在矛盾而无人察觉；它可以由推理器自动检验。
\textbf{完备性}针对「该答的答不了」——衡量标准不是包罗万象，
而是覆盖预先列出的业务问题清单，第 3 章的能力问题正是这份清单的制度化形式。
\textbf{简约性}针对「过度建模」——与任务无关的概念只会增加维护成本与共识成本。
\textbf{可扩展性}针对「一改就塌」——新概念加入不应推翻既有结构。
\textbf{共享性}针对「无人使用」——专家不认、系统不接的本体是失败的本体，无论多精致。
注意五条判据全部\textbf{可操作}：每一条都能变成评审清单上的检查项。
这与「无法收敛」的哲学追问形成了完整对照——现在可以把两种本体论并排放进一张表里了：

\begin{center}
\begin{tabularx}{\linewidth}{@{}lXX@{}}
\toprule
\textbf{维度} & \textbf{哲学本体论} & \textbf{工程本体论} \\
\midrule
研究对象 & 存在的本质、世界的本源 & 领域概念、关系与规则的形式化表示 \\
方法 & 思辨、追问 & 建模、形式化、计算 \\
判据 & 终极问题无法收敛（高熵） & 一致性可检验、可推理（低熵） \\
产出 & 哲学体系 & 可被机器处理的知识结构 \\
\bottomrule
\end{tabularx}
\end{center}

本章开头的问题至此有了答案：车间的麻烦之所以与「本体论」相关，
是因为工程本体论恰好是一门「把概念约定显式化、形式化并形成共识」的学科，
而车间三个场景的病灶——隐性、分散、歧义——正是缺少这种约定的症状。
下一节把这个对应关系整理成一个统一的视角。

\section{知识熵减：理解本体工程的统一视角}

\subsection{为什么借用「熵」这个词}

熵（entropy）在热力学中度量系统的无序程度，热力学第二定律断言：
孤立系统的熵不会自发减少——房间不收拾只会越来越乱。
香农把同名概念引入信息论，用来度量不确定性：
一条消息消除的不确定性越多，携带的信息量就越大。
本书借用这个词描述企业知识的状态，先做一个诚实的声明：
\textbf{「知识熵」是一个类比，不是严格的物理量}——
我们并不打算给企业知识算出一个以比特为单位的熵值。
借用它，是因为这个类比精确捕捉了两个事实。

第一，\textbf{知识系统放任不管必然走向混乱}，一如孤立系统熵增。
命名会漂移：每上线一个新系统就多一套字段名；
文档会腐化：流程改了文档没人改；
经验会蒸发：骨干离职带走判断力；
副本会分叉：同一份台账被复制出五个版本各自演化。
没有任何企业「决定」让知识变乱，混乱是自发的、默认的方向。
第二，\textbf{维持低熵需要持续做功}。
房间要人收拾，冰箱要耗电制冷；
知识要保持有序，同样需要持续投入建模、对齐、治理的功——
本体工程就是对知识做功的系统方法。
这笔功不白做：建模时多付出的每一分严谨，
都会在之后的检索、集成、推理、审计中按次返还。
可以概括成一句工程经验：\textbf{先交建模税，后省集成税；不交建模税，天天交沟通税}。

\subsection{四步熵减}

1.1 节诊断出三个病灶（分散、隐性、歧义），再加上「知识躺在库里不会自己工作」这层静态性，
对应的治理动作恰好是四步。包内教学资产把它们压缩成四行：

\begin{codebox}
\codeline{\textcolor{cmtgray}{\#\ ==============================}}
\codeline{\textcolor{cmtgray}{\#\ 4.\ 知识熵减的视角}}
\codeline{\textcolor{cmtgray}{\#\ ==============================}}
\codeline{\textcolor{cmtgray}{\#\ 企业知识的自然状态是高熵的：}}
\codeline{\textcolor{cmtgray}{\#\ \ \ 分散在文档、数据库、员工头脑中}}
\codeline{\textcolor{cmtgray}{\#\ \ \ 同一概念多种叫法，同一名称多种含义}}
\codeline{\textcolor{cmtgray}{\#\ \ \ 规则隐含在经验里，无法被系统校验}}
\codeblank{}
\codeline{\textcolor{cmtgray}{\#\ 本体工程的本质是知识熵减：}}
\codeline{\textcolor{cmtgray}{\#\ \ \ 隐性\ \(\rightarrow\)\ 显式：把头脑中的概念写成形式化定义}}
\codeline{\textcolor{cmtgray}{\#\ \ \ 分散\ \(\rightarrow\)\ 统一：建立全企业共享的概念词汇表}}
\codeline{\textcolor{cmtgray}{\#\ \ \ 模糊\ \(\rightarrow\)\ 精确：用公理排除非法状态}}
\codeline{\textcolor{cmtgray}{\#\ \ \ 静态\ \(\rightarrow\)\ 可推理：让机器从已有知识推导新知识}}
\codeblank{}
\codeline{\textcolor{cmtgray}{\#\ 熵减的产出：}}
\codeline{\textcolor{cmtgray}{\#\ \ \ 形式化本体（TBox）+\ 实例数据（ABox）}}
\codeline{\textcolor{cmtgray}{\#\ \ \ =\ 可查询、可推理、可复用、可解释的知识结构}}
\codeblank{}
\end{codebox}

\assetsource{ch01}{philosophy-to-engineering}

下面把四步各自展开，并各配一个车间例子。

\textbf{第一步，显性化（隐性 \(\rightarrow\) 显式）}。
把存在于头脑中的概念与规则逼问出来、写下来。
对象是场景二里王工那类「说不清」的规则型经验：
「这批料硬度偏高要降转速」经过三轮访谈逼问，变成
「当工件材料硬度超过某阈值且使用硬质合金刀具时，切削线速度上限下调百分之二十」——
条件、对象、动作、幅度全部落字。
逼问的技术（怎么提问、怎么验证、怎么处理专家之间的不一致）是第 3 章知识获取的内容。
显性化之后，知识仍然是文本，机器还读不懂，
但它已经完成了关键一跃：从「在人脑里」变成「在纸面上」。

\textbf{第二步，统一化（分散 \(\rightarrow\) 统一）}。
为全企业建立一套共享的概念词汇：声明 ERP 的「资产」、MES 的「设备」、
车间系统的「机台」是同一个概念 \texttt{Equipment}，
并为每台实际设备建立唯一标识，把三套编号映射到它。
场景三里老郑拼不出来的那条证据链，断点恰恰在这里——统一化是打通追溯链的前提。
统一化动的是各部门的「语言主权」，
因此它一半靠技术（映射表与对齐算法，第 7 章），
一半靠上一节说的社会过程（评审与共识）。

\textbf{第三步，精确化（模糊 \(\rightarrow\) 精确）}。
把含糊的表述收紧成可判定的约束。
「设备要定期保养」精确化为「每台数控机床关联一份维护计划，
计划上的保养间隔不超过其型号规定的上限」；
「运行中的设备不能维护」精确化为公理
Running \(\sqcap\) Maintenance \(\sqsubseteq\) \(\bot\)，
使「某设备同时处于运行与维护状态」成为机器可识别的非法状态。
建模实践中有一份值得随身携带的「模糊词黑名单」：
适当、酌情、及时、原则上、一般来说、相关部门——
它们每出现一次，就标记一个待精确化的点。

\textbf{第四步，可推理化（静态 \(\rightarrow\) 可推理）}。
前三步得到一份精确、统一、显式的知识资产；第四步让它「自己工作」。
形式化之后，机器可以从已有知识\textbf{推导}新知识：
由「\texttt{Lathe\_003} 是数控车床」与「数控车床是设备」自动得出「\texttt{Lathe\_003} 是设备」，
不必有人显式录入这条事实；
由钛合金加工的能力条件与 3 号车床的配置清单，自动推出「3 号车床不能承接此工单」，
顶班的小李即使不认识任何老师傅，也不会再排错。
更重要的是反向的把关：当有人试图录入与公理矛盾的断言时，
一致性检验会当场拒绝——知识库自己守卫自己的质量。
推理的机制（什么算法、什么复杂度、什么边界）是第 5 章的主题。

\subsection{熵减的产出形态：TBox 与 ABox}

四步走完，产出是什么？知识工程的标准答案是一对术语：

\begin{noteblock}
\textbf{TBox}（terminological box，术语公理集）：类、属性与公理的集合，
描述领域的概念结构——「数控车床是设备」「设备与物料不相交」「每台设备有序列号」。

\textbf{ABox}（assertional box，断言集）：个体及其事实的集合——
「\texttt{Lathe\_003} 是数控车床」「\texttt{Lathe\_003} 的功率是 15.5 千瓦」。
\end{noteblock}

一个更准确的工程比喻：TBox 是语义计算规范，ABox 是带来源的输入断言；
推理器是语义计算器，把规范应用于输入，产出可追溯的派生结论或矛盾警报。
它不裁定输入是否为现实事实，也不替有权人作风险、合规或动作决定。
两者合起来构成\textbf{知识库}（knowledge base）。
TBox 变化慢，应当由严格的治理流程控制修改；ABox 随业务流动，每天都在增删。
当 ABox 的规模增长到数亿条事实、需要专门的图存储、批量构建流水线与质量校验体系时，
这套「TBox 加超大 ABox」的工程形态就有了一个更流行的名字——知识图谱，
第 7 章专门讨论它。

工程团队常问：TBox 与 ABox 的边界划在哪里？
一条实用的判别法是看\textbf{变更频率与审批层级}：
「数控车床是设备」十年不变，改它要开评审会——放 TBox；
「3 号车床今天在维护」明天就变，随业务系统自动更新——放 ABox。
边界划错的代价是真实的：把易变事实写进 TBox，
本体版本会被业务数据搅得天天发布；
把概念约定散落在 ABox，统一化的成果又会慢慢蒸发。

于是全书的逻辑可以用熵减视角一句话串起来：
第 2 章提供描述低熵知识的数学语言，第 3 章是显性化与统一化的方法论，
第 4 章是形式化的工业标准语言，第 5 章是可推理化的机制，
第 6 章看四个行业的熵减案例，第 7 章把熵减规模化，
第 8 章让低熵知识为大模型供能，第 9 章完整跑一遍全流程。

\section{本体与相邻概念辨析}

第一次接触本体的工程师几乎都会问同一类问题：
「这不就是数据库表结构吗？」「跟 UML 类图什么区别？」「知识图谱不就是这个吗？」
这些问题问得好——本体确实与一批既有工具长得很像，
它们都在做「结构化地描述世界」这件事。
但相似只在表层，目的、语义与能力差别很大；
分不清它们，轻则选错工具，重则把项目带进沟里。
本节逐一辨析五个最常被混淆的概念。

\subsection{数据库 Schema：为存储而生，不为共享而生}

相似之处显而易见：关系数据库的 Schema 也定义「有哪些实体、什么属性、什么关系」，
实体关系图（ER 图）与本体的类与属性结构看起来是近亲。
但有四点本质差异。
\textbf{目的不同}：Schema 服务于单一应用的数据存储与完整性，
设计时优先考虑查询性能与范式，概念清晰度可以让位——拆表、冗余、代理键都是常规操作；
本体服务于跨系统的语义共享，概念保真是第一优先级。
\textbf{语义封闭}：列名 \texttt{dev\_no} 的业务含义存在于文档与数据库管理员的头脑中，
数据库本身只知道它是一个变长字符串——机器无法凭 Schema 判断
它与另一个系统的 \texttt{asset\_id} 是否同义；
本体的每个概念都携带机器可读的定义、关系与公理。
\textbf{世界假设不同}：数据库默认封闭世界——表里查不到，就当不存在；
本体默认开放世界——没有记录只说明「未知」，不等于「为假」。
第 2 章将详细讨论这条分水岭。
\textbf{推理能力不同}：数据库不会从「数控车床是设备」推出任何东西，外键也不传递语义。
常见的混用错误是把 ER 图改个画法直接宣布为「企业本体」——
得到的只是换了文件格式的私有 Schema，四个关键词一个都不满足。

\subsection{分类法 Taxonomy：只有骨架的本体}

分类法是按上下位关系组织的概念层次——设备分类树、行业分类目录、生物分类都属此类。
它与本体的关系最简单：\textbf{分类法是本体的真子集}。
本体的类层次部分就是一个分类法；
但本体在层次之外还有横向关系（设备「位于」车间、「能加工」材料）、
属性约束（功率取正数）与公理（设备与物料不相交）。
只有上下位骨架、没有血肉的分类树，回答得了「它是哪类」，
回答不了「它能干什么、和谁有关、什么状态非法」。
实务建议：企业现成的设备分类树是建本体的良好起点，但只是起点；
直接把分类树改名叫本体，是第二种常见混用错误。

\subsection{叙词表 Thesaurus：词的网络，不是概念的逻辑}

叙词表来自图书情报领域，用于检索系统的词汇控制，
典型关系有四种：用代（指向规范词）、广义词、狭义词、相关词。
它处理同义异名的能力正是统一化所需要的——这是它与本体的交集。
但叙词表的关系语义很弱：「相关词」不说明到底是什么关系；
「广义词」混装了「是一种」与「是其部分」两种完全不同的关系——
「发动机」的广义词常被标成「汽车」，可发动机并不是一种汽车，而是汽车的部件。
这种含混对人工检索无伤大雅，对机器推理是致命的。
本体则要求每种关系有明确的逻辑语义，「是一种」与「是其部分」严格分开。
顺带一提，W3C 为叙词表类资源提供了专门的表示语言 SKOS，
适合「只需要词汇控制、不需要重逻辑」的场合——这是合理的轻量选择，
但选它的人要清楚：自己建的是词表，不是本体。

\subsection{UML 类图：设计系统的图，不是描述世界的理论}

UML 类图与本体的亲缘关系最近：类、属性、关联、继承一应俱全，
受过面向对象训练的工程师上手本体建模确实更快。
差异在三个层面。
\textbf{用途}：类图描述「我们将要构建的软件系统长什么样」，是设计期工件，
系统交付后往往束之高阁；
本体描述「领域世界里什么为真」，是运行期知识，系统在线时被持续查询与推理。
\textbf{语义}：UML 的关联与聚合没有统一的形式语义，同一张图两个人可以读出两种意思，
也没有通用的推理器能对类图做一致性判定；
本体语言建立在描述逻辑之上，每条公理的含义由模型论严格定义，
推理器可以机械地判定「这套定义有没有矛盾」「这个类可不可能有实例」。
\textbf{边界}：类图属于单个系统；本体生来就是给多个系统共享的。
第三种常见混用错误由此而来：用画图工具「画」了一个本体，
却从不运行推理器、从不被任何系统查询——形式化的学费交了，形式化的红利一分没领。

\subsection{知识图谱：本体的规模化兄弟}

知识图谱（knowledge graph）与本体的关系最密切，也最常被混为一谈。
直观地说，知识图谱是用图结构组织的大规模事实集合：
节点是实体，边是关系——「\texttt{Lathe\_003} 位于二号车间」「\texttt{Lathe\_003} 类型是数控车床」。
用本章已有的术语，可以给出一个工程上够用的近似：
\textbf{知识图谱约等于本体（TBox，概念层）加上大规模实例（ABox，数据层）}。
二者侧重不同：本体研究的重心在概念与公理的质量——定义是否严谨、约束是否自洽；
知识图谱的重心在实例的规模与覆盖——亿级三元组的抽取、存储、查询与质量。
没有本体的图谱是「语义不明的边的堆积」：
属性名随抽取程序漂移，同一关系十种写法，越大越乱；
没有图谱的本体是「没有事实的法律条文」：公理精美，无事可断。
健康的工程形态是\textbf{本体驱动图谱}：先有概念契约，再按契约规模化生产事实——
这正是第 7 章的主线。

\subsection{一张对照表与一条光谱}

把五个概念放进同一张表：

\begin{center}
\begin{tabularx}{\linewidth}{@{}lXX@{}}
\toprule
\textbf{概念} & \textbf{与本体的相似之处} & \textbf{与本体的关键差异} \\
\midrule
数据库 Schema & 定义实体、属性与关系 & 服务单一应用；语义在文档与人脑中；封闭世界；无推理 \\
分类法 & 与本体的类层次同构 & 只有上下位关系，缺少横向关系、属性约束与公理 \\
叙词表 & 处理同义词与词汇控制 & 词间关系语义弱，混淆「是一种」与「是其部分」，不可逻辑推理 \\
UML 类图 & 类、属性、关联、继承俱全 & 设计期工件，无统一形式语义与推理器，限于单系统 \\
知识图谱 & 共享同一套概念层（TBox） & 重心在大规模实例层；缺本体约束的图谱会语义漂移 \\
\bottomrule
\end{tabularx}
\end{center}

也可以把它们排在一条\textbf{表达力光谱}上：
受控词表、分类法、叙词表、概念模型（ER 图与 UML）、轻量级本体、重量级本体——
从左到右，语义约束逐级增强，机器能做的事越来越多，建模与维护成本也同步上升。
工程选型的原则不是「越右越好」，而是\textbf{够用即可}：
只做检索与导航，词表可能就够；要做一致性把关与自动分类，就得走到逻辑本体这一端。
第 4 章介绍语言家族时会再次回到这条光谱，并给出选型建议。

最后把本节的四个「不要」整理成一份避坑清单：
不要把 ER 图换个文件格式就叫本体——缺共享语义，且把封闭世界假设偷偷带了进来；
不要把分类树当成完整本体——缺横向关系与公理；
不要用画图代替形式化——不运行推理器，就永远拿不到形式化的红利；
不要把「知识图谱」与「本体」混着用——一个重实例规模，一个重概念质量。
下次在评审会上听到这些说法时，你知道该怎么接话了。

\section{AI 时代的定位：给大模型一个受控语义权威层}

前几节的论证都还站在「传统」信息化的地基上。
近几年，大语言模型（LLM，Large Language Model）把「机器理解自然语言」从科幻拉进了日常，
一个自然的疑问随之而来：大模型什么都能聊，
本体这种「笨重」的符号工程还有必要吗？
本节的回答是：不但有必要，
而且\textbf{大模型的工程化落地恰恰让本体从可选项变成了必选项}。
理由要从大模型的软肋说起。

\subsection{幻觉：统计生成的内生缺陷}

LLM 的工作原理可以概括为大规模的统计推断：
模型在海量文本上学习语言的条件概率分布，
生成时根据上文预测「最可能出现的下一个词」，逐词续写。
这个机制决定了它优化的目标是\textbf{语言上的似然}，而不是\textbf{事实上的为真}。
模型参数里压缩存储的是语言模式，而不是一条条带出处、可核查的事实；
当被问到训练数据中缺失、稀疏或互相矛盾的内容时，
它仍然会按概率续写出\textbf{看起来通顺合理}的文本——这就是\textbf{幻觉}（hallucination）。
请注意：幻觉不是偶发故障，而是这种生成机制的\textbf{内生属性}；
更大的模型与更好的对齐训练能显著降低幻觉的频率，
但没有任何已知方法能把它压到零。
教学资产的第一节给出了简明的概括：

\begin{codebox}
\codeline{\textcolor{cmtgray}{\#\ ==============================}}
\codeline{\textcolor{cmtgray}{\#\ 1.\ 大模型的知识困境}}
\codeline{\textcolor{cmtgray}{\#\ ==============================}}
\codeline{\textcolor{cmtgray}{\#\ LLM\ 的本质是统计推断：根据上下文预测最可能的下一个词}}
\codeline{\textcolor{cmtgray}{\#\ \ \ 它"知道"的是语言模式，而不是经过核实的事实}}
\codeline{\textcolor{cmtgray}{\#\ \ \ 当训练数据缺失或模糊时，它会生成"看起来合理"的内容——幻觉}}
\codeblank{}
\codeline{\textcolor{cmtgray}{\#\ 幻觉在工程场景中的代价：}}
\codeline{\textcolor{cmtgray}{\#\ \ \ 设备参数答错\ \(\rightarrow\)\ 工艺事故}}
\codeline{\textcolor{cmtgray}{\#\ \ \ 法规条款编造\ \(\rightarrow\)\ 合规风险}}
\codeline{\textcolor{cmtgray}{\#\ \ \ 维修步骤臆造\ \(\rightarrow\)\ 安全事故}}
\codeblank{}
\end{codebox}
\color{black}

\assetsource{ch01}{ontology-in-ai-era}

注释列出的三种代价值得逐一掂量：
设备参数答错，轻则废件、重则伤人；
法规条款编造，企业可能在审计或诉讼中付出真金白银；
维修步骤臆造，等于把安全规程交给一台「自信的复读机」。
闲聊场景下幻觉是趣闻，工程场景下幻觉是责任事故。
有人寄希望于检索增强生成（RAG，Retrieval-Augmented Generation）——
先检索相关文档片段，再让模型基于片段作答。
RAG 确实显著缓解幻觉，但没有根治：
检索回来的仍是非结构化文本，新旧版本可能冲突，片段拼接可能断章取义，
而文本与文本之间的矛盾，机器在文本层面发现不了。
要害在于：\textbf{系统里缺少可由机器核对的语义契约、受控事实来源与证据链}。

\subsection{三种权威分离：语义、事实与决定}

“单一事实源”常被说成一句口号，但工程上必须拆成三种责任。
\textbf{语义权威}回答词汇、关系、约束、查询与推理 profile 是什么；本书由
Semantica 的内容绑定 package 承担。\textbf{事实权威}回答某台设备在某时刻的功率、
状态和来源；它来自受控台账、测量或签名记录，Semantica 不凭空创造。
\textbf{决策权威}回答谁能批准发布、写入或设备动作；语义符合不能自动授予权限。

因此中心论点应写得更严格：\textbf{工程级 AI 的每条事实性输出，都应追溯到受控事实，
并在具名 Semantica package/snapshot 的语义下被查询或校验；每个动作还要经过独立授权}。
TBox 是概念契约，ABox 是带来源和时效的事实载体，receipt 则绑定一次具体复算。
解法不是“更努力训练模型”，而是改变架构和责任边界。

\subsection{两条回答路径的对照实验}

架构上的差异落到一次具体问答上是什么样子？
设问题为：「\texttt{Lathe\_003} 的最大功率是多少？可以加工钛合金吗？」
教学资产给出了两条路径的逐步对照：

\begin{codebox}
\codeline{\textcolor{cmtgray}{\#\ 路径A\ ——\ LLM\ 直接回答（无本体）：}}
\codeline{\textcolor{cmtgray}{\#\ \ \ 模型基于训练语料中"车床通常……"的统计模式作答}}
\codeline{\textcolor{cmtgray}{\#\ \ \ \(\rightarrow\)\ 数值可能凭空生成，能力判断无依据，错误无法被发现}}
\codeblank{}
\codeline{\textcolor{cmtgray}{\#\ 路径B\ ——\ 本体约束下回答（有本体）：}}
\codeline{\textcolor{cmtgray}{\#\ \ \ Step\ 1\ \ 实体链接：'Lathe\_003'\ \(\rightarrow\)\ mfg:Lathe\_003（本体中的个体）}}
\codeline{\textcolor{cmtgray}{\#\ \ \ Step\ 2\ \ SPARQL查询事实：}}
\codeline{\textcolor{cmtgray}{\#\ \ \ \ \ \ \ \ \ \ \ SELECT\ ?p\ WHERE\ \{\ mfg:Lathe\_003\ mfg:hasPower\ ?p\ \}}}
\codeline{\textcolor{cmtgray}{\#\ \ \ \ \ \ \ \ \ \ \ \(\rightarrow\)\ 15.5\ kW（来自设备台账，有出处）}}
\codeline{\textcolor{cmtgray}{\#\ \ \ Step\ 3\ \ 能力校验：}}
\codeline{\textcolor{cmtgray}{\#\ \ \ \ \ \ \ \ \ \ \ ASK\ \{\ mfg:Lathe\_003\ mfg:canProcess\ mfg:TitaniumAlloy\ \}}}
\codeline{\textcolor{cmtgray}{\#\ \ \ \ \ \ \ \ \ \ \ \(\rightarrow\)\ false（本体中无此断言）}}
\codeline{\textcolor{cmtgray}{\#\ \ \ Step\ 4\ \ 生成回答：功率\ 15.5\ kW（来源：设备台账\ 2026-03）；}}
\codeline{\textcolor{cmtgray}{\#\ \ \ \ \ \ \ \ \ \ \ 当前知识库中没有该设备可加工钛合金的记录（区别于"不能"）}}
\codeblank{}
\codeline{\textcolor{cmtgray}{\#\ 关键差异：}}
\codeline{\textcolor{cmtgray}{\#\ \ \ 路径B的每个事实都可溯源到结构化知识；}}
\codeline{\textcolor{cmtgray}{\#\ \ \ "知识库中没有记录"与"事实为否"被严格区分（开放世界假设，见第2章）}}
\end{codebox}

\assetsource{ch01}{ontology-in-ai-era}

路径 A 一步到位：模型凭训练语料中「车床通常十几千瓦」的统计印象直接作答，
流畅、自信、可能全错，而且\textbf{错了也没有任何环节能发现}。
路径 B 走了四步，每一步都值得停下来看。
\textbf{第一步，实体链接}：把自然语言中的「Lathe\_003」锚定到本体中的个体
\texttt{mfg:Lathe\_003}，从源头杜绝张冠李戴——问的是 3 号车床，答的却是 13 号。
\textbf{第二步，结构化查询}：功率值用 SPARQL（第 4 章）从知识库中\textbf{查}出来，
而不是从模型参数里\textbf{回忆}出来；15.5 千瓦带着出处（设备台账）与入库时间。
\textbf{第三步，能力校验}：「能否加工钛合金」转成 ASK 查询，交给
Semantica 对指定 snapshot 执行；返回 false 只表示该图模式没有匹配证据，
不等于逻辑上证明“不能”。
\textbf{第四步，受约束生成}：模型把结构化结果组织成自然语言，并如实陈述知识边界。

注意路径 B 的最后一步：回答是「知识库中没有该设备可加工钛合金的记录」，
而不是「该设备不能加工钛合金」。
\textbf{「无记录」与「不能」被严格区分}——这来自本体的开放世界假设（第 2 章），
是 AI 产品可信度的关键细节：前者提示人去补充核实，后者直接把门关死。
两条路径的本质区别可以一言以蔽之：
路径 A 的错误不可发现，路径 B 的每一步都\textbf{可审计}。
工程系统要的从来不是「显得聪明」，而是\textbf{可问责}。

\subsection{三类硬约束}

把路径 B 的把关逻辑一般化，本体为 AI 系统提供三类硬约束：

\begin{codebox}
\codeline{\textcolor{cmtgray}{\#\ ==============================}}
\codeline{\textcolor{cmtgray}{\#\ 3.\ 本体作为幻觉控制层}}
\codeline{\textcolor{cmtgray}{\#\ ==============================}}
\codeline{\textcolor{cmtgray}{\#\ 本体为\ AI\ 系统提供三类硬约束：}}
\codeblank{}
\codeline{\textcolor{cmtgray}{\#\ (1)\ 词汇约束\ ——\ AI\ 只能使用本体中定义的类与属性}}
\codeline{\textcolor{cmtgray}{\#\ \ \ \ \ 防止编造不存在的概念（"超导车床"）}}
\codeblank{}
\codeline{\textcolor{cmtgray}{\#\ (2)\ 公理约束\ ——\ 违反公理的陈述直接被拒绝}}
\codeline{Running\ \(\sqcap\)\ Maintenance\ \(\sqsubseteq\)\ \(\bot\)\ \ \ \ \ \ \ \ \#\ "设备正在运行中维护"是非法陈述}
\codeline{Equipment\ \(\sqcap\)\ Material\ \(\sqsubseteq\)\ \(\bot\)\ \ \ \ \ \ \ \ \ \ \#\ "把车床作为原材料投料"是非法陈述}
\codeblank{}
\codeline{\textcolor{cmtgray}{\#\ (3)\ 范围约束\ ——\ 数值必须落在物理合理区间}}
\codeline{\textcolor{cmtgray}{\#\ \ \ \ \ 功率\ ∈\ [0,\ 10000]\ kW；温度\ \(\ge\)\ -273.15\ °C}}
\codeline{\textcolor{cmtgray}{\#\ \ \ \ \ （SHACL\ 实现见第7章质量校验）}}
\codeblank{}
\end{codebox}

\assetsource{ch01}{ontology-in-ai-era}

\textbf{词汇约束}盯住「说的词存不存在」。
AI 的输出在落库或执行之前要过一道词汇关：
提到的类、属性、个体必须能在本体的词汇签名中找到——
「超导车床」「量子热处理」这类生造概念直接打回。
看似简单，拦截效果却最立竿见影，因为编造专名正是幻觉最常见的形态。
工程实现手段包括实体链接、面向受限词表的生成、生成后校验等，第 8 章详述。

\textbf{公理约束}盯住「说的话合不合法」。
每个词都合法，组合起来仍可能违宪：
「这台设备正在运行中维护」里没有一个生造词，
但它违反公理 Running \(\sqcap\) Maintenance \(\sqsubseteq\) \(\bot\)。
做法是在隔离的候选 Dataset 上调用 manifest 声明的校验/推理 profile；
只有 runtime 实际支持并绑定 oracle 的检查才能作证。完整 DL 未被当前 Semantica
承诺时必须 fail closed，不能借“宪法法院”比喻假称已完成。

\textbf{范围约束}盯住「数值在不在理」。
词汇与逻辑都过关的陈述，数值仍可能荒谬：
功率十五万千瓦的车床、低于绝对零度的炉温、早于下单日期的交货承诺。
范围约束为数值属性声明物理与业务上的合理区间，
工程上通常用 SHACL 形状校验实现；本书的原生例证是第 7 章 ch07 包场景，
shape、正例、单因反例与 violation oracle 均由同一 runner 绑定。

三类约束分层把关，汇总如下：

\begin{center}
\begin{tabularx}{\linewidth}{@{}lXX@{}}
\toprule
\textbf{约束类型} & \textbf{拦截的典型错误} & \textbf{实现机制与详述章节} \\
\midrule
词汇约束 & 编造概念：「超导车床」 & 实体链接、受限生成与输出校验（第 8 章） \\
公理约束 & 矛盾陈述：「运行中维护」 & 推理器一致性检验（第 2、5 章） \\
范围约束 & 失实数值：「功率十五万千瓦」 & SHACL 形状校验（第 7 章） \\
\bottomrule
\end{tabularx}
\end{center}

\subsection{互补而非替代}

行文至此要防止一个反向的误解：本体不是来取代大模型的，大模型也淘汰不了本体。
两者的强项几乎完美互补：

\begin{codebox}
\codeline{\textcolor{cmtgray}{\#\ ==============================}}
\codeline{\textcolor{cmtgray}{\#\ 4.\ LLM\ 与本体的互补结构}}
\codeline{\textcolor{cmtgray}{\#\ ==============================}}
\codeline{\textcolor{cmtgray}{\#\ \ \ LLM\ 擅长：自然语言理解、意图识别、文本生成、模糊匹配}}
\codeline{\textcolor{cmtgray}{\#\ \ \ 本体擅长：事实存储、一致性检验、逻辑推理、可解释溯源}}
\codeblank{}
\codeline{\textcolor{cmtgray}{\#\ \ \ 分工模式：}}
\codeline{\textcolor{cmtgray}{\#\ \ \ 用户自然语言\ \(\rightarrow\)\ [LLM]\ 解析意图、生成查询}}
\codeline{\textcolor{cmtgray}{\#\ \ \ \ \ \ \ \ \ \ \ \ \ \ \ \(\rightarrow\)\ [本体/知识图谱]\ 检索事实、校验约束、执行推理}}
\codeline{\textcolor{cmtgray}{\#\ \ \ \ \ \ \ \ \ \ \ \ \ \ \ \(\rightarrow\)\ [LLM]\ 把结构化结果组织成自然语言回答}}
\codeline{\textcolor{cmtgray}{\#\ \ \ （实现模式详见第8章：GraphRAG、Text2SPARQL、本体引导Agent）}}
\codeblank{}
\end{codebox}

\assetsource{ch01}{ontology-in-ai-era}

分工模式已经相当清晰：LLM 站在人机界面一侧，
负责听懂意图、把问题翻译成结构化查询、把结构化结果还原成自然语言——
它是出色的\textbf{翻译官与文案}；
Semantica 站在语义执行与证据一侧，负责 package、查询、校验、受限推理、版本和 receipt；
现实事实仍由具名来源负责。更准确的分工是“神经网络管理解与表达，受控来源管事实，
Semantica 管可执行语义，具名主体管决定”。这属于神经符号（neuro-symbolic）路线；
工业界的具体实现模式——GraphRAG、Text2SPARQL、本体引导的智能体（Agent）——
将在第 8 章逐一拆解并配上代码。
顺带说一句，这种互补是双向的：本体工程也在受益于大模型——
用 LLM 辅助抽取概念、起草公理、生成跨系统映射，
能显著降低第 3 章那些方法论步骤的人力成本。
一句话总结本节：\textbf{大模型负责把话说漂亮，本体负责把事说对}。

\section{产业图景}

本体不是书斋技术。过去二十多年，它沉淀出一套国际标准、若干公共知识资源
和几条清晰的产业落地模式。本节速览这张图景——
不求全景，但求读者知道自己学的东西在产业坐标系中的位置。

\subsection{W3C 语义网标准生态}

2001 年，万维网发明人 Tim Berners-Lee 与合作者在《科学美国人》上发表《语义网》一文，
描绘了「让网页内容对机器有意义」的愿景。
二十多年过去，「整个 Web 变成一个大知识库」的原始愿景并未按字面实现，
但这场运动在 W3C（万维网联盟）沉淀出一套层层叠放、至今活跃的技术栈，
成为本体工程事实上的通用语言：

\begin{center}
\begin{tabularx}{\linewidth}{@{}lXX@{}}
\toprule
\textbf{标准或规范} & \textbf{解决的问题} & \textbf{本书章节} \\
\midrule
RDF & 用「主语——谓语——宾语」三元组表示一切事实的数据模型 & 第 4 章 \\
RDFS & 在 RDF 之上声明类、属性与简单层次的轻量词汇 & 第 4 章 \\
OWL & 完整的本体语言，以描述逻辑为根基，支持自动推理 & 第 4 章、第 5 章 \\
SPARQL & 面向三元组库的标准查询语言 & 第 4 章 \\
SWRL & 在 OWL 之上书写业务规则 & 第 5 章 \\
SHACL & 声明数据「形状」，批量校验实例质量 & 第 7 章 \\
\bottomrule
\end{tabularx}
\end{center}

表中各项除 SWRL 是成员提交规范外，均为 W3C 推荐标准。
标准化的意义怎么强调都不过分：
概念与事实一旦写成 RDF 与 OWL，就获得了\textbf{厂商中立的可移植性}——
建模工具（如 Protégé）、编程框架（如 Jena）、各家图数据库都讲同一种语言，
知识写一次，换存储、换工具不必重写。
对比各自为政的私有 Schema 世界，这正是「共享」这个关键词在产业尺度上的兑现。
对企业读者，标准栈还有一层现实意义：人才与知识资产的可流动性——
会 SPARQL 的工程师换一家公司仍然会 SPARQL，
按标准构建的知识库也不会被某一家供应商锁死。

\subsection{公共本体：从搜索引擎到生物医学}

标准之上生长出一批公共知识资源，举几个不同领域的代表。
面向 Web 的 schema.org 由主流搜索引擎厂商于 2011 年联合发起，
为网页提供结构化标注词汇——今天你搜一道菜谱能看到带评分与用时的卡片，
背后就是海量网页嵌入的 schema.org 标注；它是「轻量本体加海量采用」路线的标杆。
生物医学是重量级本体最成功的领域：
基因本体（Gene Ontology）是基因功能注释的事实标准，
支撑着全球基因组数据的可比与共享；
临床术语系统 SNOMED CT 覆盖数十万医学概念，被多个国家的医疗信息系统采用。
在更抽象的层面，还有 BFO、DOLCE 等\textbf{上层本体}（upper ontology）——
它们不描述具体领域，而为「物体」「过程」「品质」这类跨领域范畴提供统一框架，
可以视作亚里士多德十范畴事业的现代续篇，
其中 BFO 已进入国际标准体系（ISO/IEC 21838 系列）。
工业领域则有面向流程工业全生命周期数据集成的 ISO 15926，
以及工业 4.0 资产管理壳（Asset Administration Shell）对语义化描述的采纳。
这串名单传递的信息是：\textbf{本体早已不是实验室玩具，而是多个行业的在役基础设施}；
当然，各领域落地深浅不一，也不必神化。

\subsection{数据平台的「本体层」模式}

近年企业数据架构中出现了一个值得注意的模式：
在数据湖与数据仓库之上增加一个\textbf{本体层}（ontology layer），
把业务对象（设备、订单、客户）、对象间关系乃至可执行的业务动作统一定义在本体中；
上层应用与 AI 不再直接面对几千张物理表，
而是面对几百个有定义、有关系、有权限语义的业务概念。
收益是系统性的：语义只定义一次，处处复用；
数据血缘与访问权限可以挂在业务概念而非物理列上；
AI 的输入输出有了天然的约束框架——正是上一节三类硬约束的落点。
商业世界里，已有数据平台厂商（较知名的如 Palantir 的 Foundry 平台）
把「本体」作为产品的核心概念公开宣传；
各大云厂商与图数据库厂商也普遍提供知识图谱构建与托管服务。
本书不评价任何商业产品的优劣，只指出一点：
这套模式的开放技术路线，恰好就是本书第 4 章（语言）、第 7 章（图谱构建）
与第 8 章（LLM 融合）合起来要教的东西。

\subsection{工业知识图谱的实践方向}

最后看与本书案例最贴近的工业场景。教学资产归纳了三个方向：

\begin{codebox}
\codeline{\textcolor{cmtgray}{\#\ ==============================}}
\codeline{\textcolor{cmtgray}{\#\ 5.\ 工业界实践方向}}
\codeline{\textcolor{cmtgray}{\#\ ==============================}}
\codeline{\textcolor{cmtgray}{\#\ 数据分析平台：以本体层（Ontology\ Layer）作为企业数据的语义中枢，}}
\codeline{\textcolor{cmtgray}{\#\ \ \ 业务对象、关系、行动在本体中统一定义，AI\ 在其上运行}}
\codeline{\textcolor{cmtgray}{\#\ 制造企业：设备、工艺、质量知识图谱，支撑故障诊断与工艺推荐}}
\codeline{\textcolor{cmtgray}{\#\ 监管行业：法规条款本体化，让合规判断可溯源、可审计}}
\codeblank{}
\codeline{\textcolor{cmtgray}{\#\ 共同点：让\ AI\ 的每个结论都能追溯到一条结构化的知识定义，}}
\codeline{\textcolor{cmtgray}{\#\ 从根源上抑制大模型幻觉——这正是本书后续各章要构建的能力}}
\end{codebox}
\color{black}

\assetsource{ch01}{ontology-in-ai-era}

在制造业内部，公开案例较多的应用集中在几条线上。
\textbf{设备故障诊断}：把「故障模式、可能原因、排查步骤、处置措施」的链条本体化，
维修工单与诊断系统共享同一套故障概念，新故障可以沿着链条做归因检索。
\textbf{工艺知识复用}：把材料、工艺、参数、质量结果之间的依赖关系图谱化，
新订单进来时检索相似工艺并给出可解释的依据——正是场景二里王工经验的归宿。
\textbf{供应链与合规}：多级供应商关系、物料合规属性入图，
支撑风险传导分析与审计追溯——正是场景三里老郑需要的那条证据链。
强监管行业（航空、核电、制药）的动力更直接：
法规条款本体化之后，合规判断从「人查文件」变成「机器查图谱加人复核」，
每一步可溯源、可审计。
对这些方向，本书的态度是客观稳妥：它们都有公开的实践与标准化活动支撑，
但成熟度参差，不存在「上了图谱就智能」的捷径——
第 6 章将用制造调度、自动驾驶、建筑信息模型（BIM）、航空故障模式分析（FMEA）
四个案例展示「可行的做法」长什么样。
读完本书，你在这张产业图景中的位置也就清楚了：
成为\textbf{具备「本体层」设计与实施能力的人}——
无论职位名片上印的是本体工程师、知识工程师，还是语义架构师。

\section{全书结构与三条学习路径}

\subsection{九章的任务分解}

全书九章构成「理论 \(\rightarrow\) 方法 \(\rightarrow\) 语言 \(\rightarrow\) 推理
\(\rightarrow\) 应用 \(\rightarrow\) 规模化 \(\rightarrow\) AI 融合 \(\rightarrow\) 实战」的链路，
每一章解决一个明确的问题：

\begin{itemize}
\item \textbf{第 2 章（基础理论）}：本体由哪些构件组成——类、属性、公理、实例；
      背后的数学语言是什么——描述逻辑与一阶逻辑；
      开放世界假设与单调性这些「世界观」级别的设定如何影响一切后续设计；
\item \textbf{第 3 章（构建方法论）}：从零开始建本体的工序——能力问题（CQ）定范围，
      Ontology 101 七步法走流程，METHONTOLOGY 管生命周期，OntoClean 查层次质量；
\item \textbf{第 4 章（语言与工具）}：把设计落成代码——RDF、RDFS、OWL 语言家族，
      SPARQL 查询，Protégé 建模环境；
\item \textbf{第 5 章（推理机制）}：机器怎么「想」——Tableau 算法、SWRL 规则、
      时态与概率扩展，以及推理的复杂度边界；
\item \textbf{第 6 章（应用案例）}：制造调度、自动驾驶、BIM、航空 FMEA
      四个领域的建模实战；
\item \textbf{第 7 章（知识图谱）}：当实例规模到亿级——抽取、对齐、融合、
      存储选型与 SHACL 质量门禁；
\item \textbf{第 8 章（本体与大模型）}：幻觉控制架构、GraphRAG、Text2SPARQL、
      本体引导 Agent；
\item \textbf{第 9 章（综合实战）}：把制造业 CQ、查询、规则、夹具与 oracle
      绑定到 Semantica 单一 package/runner；保留旧双栈片段作实现史比较，同时诚实暴露缺口。
\end{itemize}

Semantica 中九个 package 与九章一一对应：
\texttt{semantica.chapter\_packages.vol1.ch01} 至
\texttt{semantica.chapter\_packages.vol1.ch09}。这棵树的枝条不再指向书旁的
迁移前的配套目录，而指向每章的 manifest、asset、scenario 与 gap；
理论地基仍在书中，机器可执行语义则只有一个落点。

\subsection{章节依赖与最短路径}

各章不是均匀的线性链。主干是第 2 章到第 5 章的递进——理论、方法、语言、推理层层铺垫；
第 6 章的案例同时取材于第 3、4、5 章；
第 7、8 两章并列地立在第 4 到第 6 章之上，相互独立、可以换序阅读；
第 9 章是总收口，对前八章均有回指。
赶时间的读者可以走\textbf{最短可用路径}：
第 4 章（会读写语言）\(\rightarrow\) 第 7 章（理解 Dataset 与质量门禁）
\(\rightarrow\) 第 9 章（复算具名场景）——三章读完能掌握最小可验证链路。
但要诚实提示代价：跳过第 2 章的人，
在系统第一次出现「推理结果反直觉」时会束手无策；
跳过第 3 章的人，建出的本体往往在第三个月开始返工。
最短路径适合「先跑起来再回头补课」的工作风格——前提是别省略补课环节。

\subsection{三条学习路径}

结合常见的读者画像，给出三条路径建议，
与包内 \texttt{book-roadmap} 教学资产第三节的阅读路径一致：

\begin{itemize}
\item \textbf{理论完整路径}（高校课程或系统自学）：按 1 至 9 章顺序精读，
      重点逐节消化第 2 章描述逻辑与第 5 章推理机制；
      适合把本书当一学期课程使用的读者，产出是完整自洽的知识体系；
\item \textbf{工程速成路径}（在职工程师、项目驱动）：1 \(\rightarrow\) 3（CQ 与七步法）
      \(\rightarrow\) 4（OWL 与 SPARQL）\(\rightarrow\) 7（建库）\(\rightarrow\) 9（实战）；
      Tableau 算法与概率推理先跳过，项目中遇到推理问题时回查第 5 章；
      产出是短周期内可以上手真实项目的动手能力；
\item \textbf{AI 应用路径}（大模型应用开发者）：1 \(\rightarrow\) 8（先看本体能为 LLM 做什么）
      \(\rightarrow\) 4（补语言基础）\(\rightarrow\) 7（GraphRAG 的图谱底座）
      \(\rightarrow\) 5（推理校验）\(\rightarrow\) 9；
      重点是第 8 章幻觉控制架构与第 7 章 SHACL 校验；
      产出是给现有 LLM 应用加装受控事实、Semantica 语义门禁与审计链的能力。
\end{itemize}

无论选哪条路径，有两条建议不变。
其一，第 9 章务必动手，但“动手”现在指查看 manifest、运行具名 scenario、
核对 exact oracle，再观察 release verification 为什么仍阻断。
其二，善用包清单：正文片段下方保留历史来源锚点，构建器则通过迁移账定位到
Semantica 的唯一 asset。第 2、3 章以概念推演和 CQ 为主；第 4 章是 RDF/OWL/SPARQL；
第 5 章同时展示可执行的受限规则与明确不支持的推理 profile；第 6 章案例仍未形式化；
第 7 章有原生 SHACL 场景；第 8 章动作合同仍阻断；第 9 章只有 CQ1 拥有 exact oracle。
红灯与绿灯同样是学习材料，不以另装第三方引擎来改写结果。

\section{本章小结与练习}

\subsection{要点回顾}

\begin{itemize}
\item 企业知识的普遍现状是\textbf{信息都在、知识不通}，三个病灶：
      分散（系统孤岛）、隐性（经验在人脑）、歧义（同名异义与同义异名）——合称高熵状态；
\item 哲学本体论两千多年的追问\textbf{无法收敛、不可计算}，
      但留下顶层分类、属加种差定义法与同一性判据三笔遗产；
\item 工程转向的标志是 Gruber 到 Studer 的定义：
      本体是\textbf{共享概念化的形式化、显式规范}；
      四个关键词各划掉一类冒牌货——没有概念化的（裸数据）、不显式的（个人经验）、
      不形式化的（自然语言文档）、不共享的（个人笔记）；
\item 小写复数的 ontologies 宣告本体是\textbf{工件}：有版本、有任务边界、可评审可重构；
      好坏以一致性、完备性、简约性、可扩展性、共享性五条工程判据衡量；
\item 本体工程即\textbf{知识熵减}四步：显性化、统一化、精确化、可推理化；
      产出 TBox（条文）加 ABox（事实），规模化之后就是知识图谱；
\item 本体不同于数据库 Schema（私有语义、封闭世界）、分类法（只有骨架）、
      叙词表（弱语义）、UML 类图（无形式语义与推理）；
      与知识图谱是概念层与实例层的关系；
\item AI 时代本体的新身份是\textbf{语义契约}：
      Semantica 以 package 化词汇、公理、范围、查询与 shape 执行已声明门禁；
      事实来自受控来源，决定来自具名主体，LLM 只负责理解与表达。
\end{itemize}

导读里许诺过一句话的版本，现在可以兑现了：
\textbf{本体是让机器「懂」业务的基础设施}——
「懂」的工程含义，是用共享的概念契约去查询、校验、推理每一条业务陈述；
读完后续八章，你将有能力把这句话落成一个边界、输入、输出与证据均可复核的系统；
对尚未实现的能力保留红灯，也是这种能力的一部分。

\subsection{思考题}

\begin{noteblock}
\textbf{思考题}
\begin{enumerate}
\item 在你熟悉的组织里找一个「信息都在、知识不通」的实例，
      指出它的病灶属于分散、隐性、歧义中的哪几种，并估计它每年造成的隐性成本。
\item 用概念化、显式、形式化、共享四个关键词逐一检验：
      你公司的《设备管理制度》文档是不是一个本体？它最缺哪个关键词？补齐需要做什么？
\item 为你所在企业的「设备同一性」设计一个工程判据，并讨论三种情形：
      大修更换主要部件、整机搬迁到另一工厂、改造升级后型号变更。
      你的判据在哪种情形下会失效？
\item 选一条你听过的老师傅经验（任何行业皆可），先做显性化（写成完整的条件句），
      再做精确化（给出可判定的阈值与范围），
      最后尝试用 \(\sqsubseteq\)、\(\sqcap\) 等符号写出它的骨架。哪一步最难？为什么？
\item 判断并说明理由：「我们已经有完整的数据库 ER 图，所以我们已经有本体了。」
      提示：从四个关键词与开放世界假设两个角度检查。
\item 词汇、公理、范围三类硬约束分别能拦截下列哪条错误输出：
      「\texttt{Lathe\_003} 的功率是一万五千千瓦」「设备 X 正在运行中维护」
      「建议采购一台超导车床」？说明每一条的拦截机制。
\item 「知识库中没有记录」与「事实为否」的区分为什么重要？
      构造一个把二者混淆会导致实际损失的业务场景，并说明系统应当如何措辞。
\item 对照本章给出的三条学习路径，为自己选定一条并写下理由、
      预计最困难的一章，以及你打算用哪个实际问题来检验学习成果。
\end{enumerate}
\end{noteblock}

\subsection{本章包资产指引}

三份教学材料的唯一机器可读正本位于 ch01 package；表中的名称是 asset id，
正文下方的旧路径只用于证明迁移来源：

\begin{center}
\begin{tabularx}{\linewidth}{@{}>{\ttfamily\footnotesize}p{0.30\linewidth}Xl@{}}
\toprule
\normalfont\textbf{Semantica asset} & \textbf{内容与阅读建议} & \textbf{关联小节} \\
\midrule
philosophy-to-engineering & 哲学与工程的逐节对照；重点看「同一个设备概念」如何被公理消解，以及五条工程判据 & 1.2--1.4 \\
ontology-in-ai-era & 两条回答路径全文、三类硬约束、LLM 与本体分工模式、产业方向 & 1.6、1.7 \\
book-roadmap & 全书结构树的来源材料；其中旧目录/双栈字样只作迁移前历史，不代表当前入口 & 1.8 \\
\bottomrule
\end{tabularx}
\end{center}

下一章我们正式走进工程本体论的内部，从四块积木——类、属性、公理、实例——开始，
把本章反复预告的那些符号（\(\sqsubseteq\)、\(\sqcap\)、\(\bot\)）变成你自己的工作语言。
